\documentclass[11pt]{article}

\usepackage[T1]{fontenc}
\usepackage{fancyvrb}

\title{ECE 154A Lab 1 Report}
\author{Shravya Salem Sathish, David Uzunov}
\date{}

\begin{document}

\maketitle

\section{Hours Spent}

We spent a total of \underline{8} hours on this lab. 

\section{\texttt{mult.s} Assembly Code}

Completed \texttt{mult.s} code pasted below with C equivalents in the comments: 

\begin{Verbatim}
add a2, zero, zero                        # a2 = 0
    LOOP:                                  
        beq a1,zero,DONE                  # while(a1!=0){ 
        andi t0,a1,1                      # t0 = a1 & 1;
        beq t0,zero,SHIFT                 # if(t0!=0){
        add a2,a2,a0                      # a2 += a0;}
        SHIFT:
            srli a1,a1,1                  # a1 = a1 >> 1;
            slli a0,a0,1                  # a0 = a2 << 1;
    j LOOP                                # } // close out the while loop
DONE:
\end{Verbatim}

\section{\texttt{div.s} Assembly Code}

Completed \texttt{div.s} code pasted below with C equivalents in the comments:

\begin{Verbatim}

# Set initial values. 
# The remainder is initially 
# equal to the dividend.
add a3, zero, a0                           # a3 = a0;

# If the divisor is zero, 
# set the quotient to 0xFF 
# and proceed to done. 
bne a1, zero, divisorNonzero               # if (a1 != 0) { divisorNonZero; }         
addi a2, zero, 0x000000FF                  # a2 = 0x000000FF;                            
j done

divisorNonzero:
# For nonzero divisors, 
# set the quotient to zero.
add a2, zero, zero                         # a2 = 0;

# If the dividend is lesser 
# than the divisor, proceed to done.
bltu a0, a1, done                          # if (a0 < a1) { done; }

# Dividend and divisor are valid, 
# so calculate their bit lengths. 
# Initially store the dividend in t0
# and the divisor in t1. Store their current
# lengths in t2 and t3, respectively.
add t0, zero, a0                           # t0 = a0;
add t1, zero, a1                           # t1 = a1;
add t2, zero, zero                         # t2 = 0;
add t3, zero, zero                         # t3 = 0;

# Right shift t0 and increment t2 
# until t0 is zero for the dividend
# bit length.
whileDividendNonzero:                      
beq t0, zero, whileDivisorNonzero         # while (t0 > 0)
srli t0, t0, 1                             # t0 >>= 1;
addi t2, t2, 1                             # t2 += 1;
j whileDividendNonzero

# Right shift t1 and increment t3 
# until t1 is zero for the divisor
# bit length.
whileDivisorNonzero: 
beq t1, zero, lenDone                     # while (t1 > 0)
srli t1, t1, 1                             # t1 >>= 1;
addi t3, t3, 1                             # t3 += 1;
j whileDivisorNonzero

lenDone:
# Left shift the divisor by the 
# difference in bit length between the
# dividend and the divisor.
sub t4, t2, t3                             # t4 = t2 - t3
sll a1, a1, t4                             # a1 <<= t4;

# Set up a for loop with t4 + 1
#  iterations for the division.
addi t4, t4, 1                             # t4 += 1;
add t5, zero, zero                         # t5 = 0;

for:                                       # for (; t5 < t4; t5++)
bge t5, t4, done

# If the remainder is greater than
# or equal to the divisor, subtract
# the divisor from the remainder, 
# left shift the quotient by 1, and
# set the quotient LSB to 1.
bltu a3, a1, else                          # if (a3 >= a1) { ... }
sub a3, a3, a1                             # a3 -= a1;
slli a2, a2, 1                             # a2 <<= 1;
ori a2, a2, 1                              # a2 |= 1;
j skip

# Else just left shift the quotient by 1,
# setting the LSB to 0.
else:                                      # else { a2 <<= 1; }
slli a2, a2, 1

# Right shift the divisor by 1
# and increment the for loop counter.
skip:
srli a1, a1, 1                             # a1 >>= 1;
addi t5, t5, 1                             # t5 += 1;
j for

done:                                      # End of division.

\end{Verbatim}

\end{document}
